\documentclass[12pt]{article}
\usepackage[T1]{fontenc}
\usepackage[utf8]{inputenc}
\usepackage{lmodern}
\usepackage{textcomp}
\usepackage{enumitem}
\usepackage{geometry}
\geometry{margin=1in}
\begin{document}
\begin{center}
Answer Key - Geography - K-12 Level Assessment 12 \
\small (Answers are ordered exactly as in the question paper.)
\end{center}

\section*{Multiple Choice}
\begin{enumerate}[label=\arabic*.]
\item \textbf{Answer:} B
\\ \textit{Options:} A) CBR - CDR 5 2.; B) Births 1 immigration 5 deaths 1 emigration.; C) Total fertility 5 3.; D) Replacement fertility is achieved.
\\ \textit{Note:} Zero population growth (ZPG) occurs when the number of births plus immigration equals the number of deaths plus emigration, resulting in a stable population size.
\item \textbf{Answer:} A
\\ \textit{Options:} A) high birth rates with high but fluctuating death rates.; B) declining birth rates with continuing high death rates.; C) low birth rates with continuing high death rates.; D) high birth rates with declining death rates.
\\ \textit{Note:} The first stage of the demographic transition is characterized by high birth rates and high but fluctuating death rates due to limited access to healthcare, poor sanitation, and high infant mortality, which leads to a stable but low population growth.
\item \textbf{Answer:} B
\\ \textit{Options:} A) urbanization.; B) gentrification.; C) suburbanization.; D) multiplier effect.
\\ \textit{Note:} Gentrification refers to the process of wealthier individuals moving into and revitalizing deteriorated urban areas, often leading to increased property values and displacement of lower-income residents.
\item \textbf{Answer:} B
\\ \textit{Options:} A) outsourcing.; B) offshoring.; C) maquiladoras.; D) locational interdependence.
\\ \textit{Note:} Offshoring refers to the practice of outsourcing business operations to a foreign third-party service provider, typically to reduce costs or access specialized skills not available domestically.
\item \textbf{Answer:} D
\\ \textit{Options:} A) population is limited by their means of subsistence.; B) all populations have the potential to increase more than the actual rate of increase.; C) wars and famine inhibit population's reproductive capacity.; D) technology's ability to raise the Earth's carrying capacity.
\\ \textit{Note:} The correct answer highlights that Malthus underestimated the impact of technological advancements in agriculture and food production, which have significantly increased the Earth's carrying capacity beyond his predictions of population growth outpacing resources.
\end{enumerate}

\section*{True / False}
\begin{enumerate}[label=\arabic*.]
\setcounter{enumi}{5}
\item \textbf{Answer:} False
\\ \textit{Options:} A) True; B) False
\\ \textit{Note:} World War I primarily influenced military technology and industrial practices, while the Second Agricultural Revolution was driven by earlier innovations in agricultural methods and tools, not directly by the war.
\item \textbf{Answer:} False
\\ \textit{Options:} A) True; B) False
\\ \textit{Note:} Taoism is considered an ethnic religion, primarily rooted in Chinese culture and philosophy, rather than a universalizing religion that seeks to attract followers from diverse cultures worldwide.
\item \textbf{Answer:} True
\\ \textit{Options:} A) True; B) False
\\ \textit{Note:} The statement is true because the United Nations frequently engages in interventionism through peacekeeping missions and diplomatic efforts to resolve conflicts and maintain international security.
\item \textbf{Answer:} False
\\ \textit{Options:} A) True; B) False
\\ \textit{Note:} Folklore encompasses diverse cultural traditions, beliefs, and practices that reflect the unique identities and histories of different groups, rather than merging them into a single, uniform tradition.
\item \textbf{Answer:} False
\\ \textit{Options:} A) True; B) False
\\ \textit{Note:} State-sponsored terrorism is not a significant issue in Canada, as the country is generally viewed as a stable and peaceful nation with a strong commitment to international law and human rights.
\end{enumerate}

\section*{Long Form Response}
\begin{enumerate}[label=\arabic*.]
\setcounter{enumi}{10}
\item \textbf{Answer:} Periphery countries primarily engage in primary economic activities, which involve the extraction and harvesting of natural resources. This includes agriculture, mining, forestry, and fishing. These activities often lead to the export of raw materials to core countries, where they are processed into finished goods. While this relationship can provide income, it often hinders the development of periphery countries by creating dependency on core countries and limiting their capacity to diversify their economies. Consequently, the focus on primary activities can lead to economic instability and a lack of technological advancement in periphery nations.
\\ \textit{Note:} correct because it identifies that periphery countries primarily engage in primary economic activities that focus on resource extraction, which leads to economic dependency on core countries and limits their development potential by constraining diversification.
\item \textbf{Answer:} Since 1980, the United States has experienced significant migration patterns primarily from Asia and Latin America. Many immigrants from Latin American countries, such as Mexico and Central America, have come seeking better economic opportunities and safety from violence, leading to vibrant Hispanic communities across the nation. Meanwhile, immigrants from Asian countries, including China, India, and the Philippines, have contributed to the tech industry and higher education sectors, enriching the U.S. economy. These migration trends have led to increased cultural diversity, influencing American cuisine, festivals, and social norms, while also presenting challenges related to integration and policy discussions. Overall, the influx of immigrants from these regions has shaped the social, economic, and cultural landscape of the United States significantly.
\\ \textit{Note:} correct because it accurately identifies the primary regions of origin for migrants-Latin America and Asia-while also addressing the social, economic, and cultural impacts of these migrations, such as the establishment of Hispanic communities and contributions to the tech industry and higher education.
\end{enumerate}

\vfill
\noindent\textit{End of Answer Key}
\end{document}